\section{One set of detecting elements for a finite family}

An exceptional point has two blocks. Its odd Selmer classes enter the
trace test in pairs, so the one-block detector of CCH3--CCH6 does not
by itself detect the new directions. We prove the original
$\mathrm{Pa.Choice}$ for any finite nonempty family of normalized
points. In particular the block sizes, the numbers of Selmer generators
and the starting coefficient fields may vary along the chain.

% BEGIN Q70C1
\paragraph{Q70C1. Integral odd Selmer generators with local witnesses.}
Take $x:\mathrm{NP}(F,n)$, $G=F.G_0$, $H=F.G'$ the kernel of the
full cyclotomic character on $G$, and $A=k[[t]]$, $E=k((t))$.
Write $r=x.\rho$, $\chi$ for the cyclotomic character in $k^\times$,
$\sigma(g)=cgc^{-1}$ and $s=-\chi(c)$. Thus $\chi|_H=1$,
$s^2=1$, and $\chi(\sigma(g))=\chi(g)$, as in CCH1.
Write $B$ for the diagonal block algebra and $V$ for the space of
off-diagonal blocks, over the indicated coefficient ring. The two
blocks have sizes $a=x.a$ and $n-a$, omitting a block of size zero.
The original power-of-$t$ span condition gives
$\operatorname{span}_E r(H)=B_E$.

Suppose $\Phi_\ell(F.\Gamma)$ holds. For a continuous $A$-linear
action of $G$ on a finite free $A$-module $M$, the module of
continuous crossed homomorphisms $G\to M$ is finite over $A$.
Here continuity means continuity modulo every $t^j$. Indeed the
kernel of the action modulo $t$ is open in $G$ and in $F.\Gamma$.
A reduced cocycle restricts there to a homomorphism to the finite
additive group $M/tM$. The hypothesis gives only finitely many such
homomorphisms, by taking $\mathbb F_\ell$-coordinates. Values on
finitely many coset representatives determine the rest. Reduction of
integral cocycles has kernel exactly $t$ times their module, by
division and continuity one level deeper. Lift a basis modulo $t$ and
successively subtract multiples of it to each precision. The resulting
coefficients converge in $A$, and equality follows pointwise by
separatedness, as in CCH7.

Apply this with conjugation on $V_A$. Inside this finite module the
original odd Selmer conditions define a submodule $S$. Indeed the
polarization and inertia conditions are linear, and local ordinary
corrections add and scale along with their cocycles. The skew
off-diagonal coboundaries form a submodule $D\subset S$ since the
action is trivial at prescribed inertia, and the defining matrix
itself gives each local correction. Since $A$ is Noetherian,
$(S/D)\otimes_A E$ is finite-dimensional. Choose an $E$-basis
represented by $\kappa_1,\ldots,\kappa_m\in S$, by extracting a
basis from the images of generators. An integral relation modulo
$D$ gives a fraction-field relation and has zero coefficients.
Conversely any class is an $E$-linear combination of these classes.
Clear all denominators and then kill the torsion in the resulting
single difference by a further power of $t$. This gives precisely
$t^b z-\sum a_j\kappa_j\in D$ with $a_j\in A$. For each
$\kappa_j$ and each $v$ take a witness from its local condition,
including the upper triangular test and the vanishing diagonal on
$F.\mu_v$. This gives all the data and fields of $k_x:\mathrm{Kap}(x)$.
For any such package its integral independence remains true against
fraction-field skew coboundaries, by clearing the coefficients and
entries of the proposed primitive simultaneously.
% END Q70C1

% BEGIN Q70C2
\paragraph{Q70C2. The quadratic trace identity for the two opposite blocks.}
Assume $\ell\ne2$ and fix any $k_x:\mathrm{Kap}(x)$ with $m$
generators. Work over $E$. The coefficient algebra of the detector
has basis $1,e_j,f_j,\eta$ ($1\le j\le m$). Among the products of
nonconstant basis vectors only $e_j f_i=f_i e_j=\delta_{ij}\eta$
can be nonzero. The first-order representation is
\[
 R(g)=\left(1+\sum_j e_j\kappa_j(g)\right)r(g).
\]
The cocycle law of $\kappa$ gives multiplicativity. Its polarization
$\kappa_j(\sigma g)=-\kappa_j(g)^T$ gives
$R(\sigma g)=\chi(g)^{1-n}R(g)^{-T}$ by the form identity for $r$.

A twisted test cocycle has precisely the shape
$Z(g)=\sum_j f_j\lambda_j(g)+\eta\phi(g)$, with $\lambda_j\in V_E$
and $\phi\in B_E$. Its equations are
\[
 Z(gu)=Z(g)+\chi(g)R(g)Z(u)R(g)^{-1},
 \qquad Z(\sigma g)=s Z(g)^T .
\]
In particular each $\lambda_j$ is an ordinary crossed homomorphism
for $\chi\operatorname{Ad}r$, with transpose sign $s$.
Suppose the $\eta$ coefficient of $\operatorname{tr}(Z(h)R(h))$
vanishes on $H$. Since products of two matrices of the shape of
$Z$ vanish, $h\mapsto(1+Z(h))R(h)$ is a representation on $H$.
Extend it linearly to $E[H]$. Write $\pi,K_j,L_j,Q$ for its
constant, $e_j$, $f_j$, and $\eta$ coefficients. Thus
$\pi(h)=r(h)$, $K_j(h)=\kappa_j(h)r(h)$ and
$L_j(h)=\lambda_j(h)r(h)$. Also $\operatorname{tr} Q=0$
on the entire group algebra. Both $K_j$ and $L_j$ satisfy
$D(uv)=D(u)\pi(v)+\pi(u)D(v)$.

Set $I=\ker\pi$. For $u,v\in I$, fix indices $i,j$ in opposite
blocks. Lift the matrix units $E_{ii},E_{jj}$ in $B_E$ to
$p_i,p_j\in E[H]$ by Q70C1. Evaluate
$\operatorname{tr}Q(u p_j v p_i)=0$. The constant terms of $u,v$
are zero, so each contributes one linear term and only the
constant terms of $p_i,p_j$ contribute. The matrix-unit trace reads
\[
 \sum_d\big(K_d(u)_{ij} L_d(v)_{ji}
              +L_d(u)_{ij} K_d(v)_{ji}\big)=0 .
\]
This identity includes the cross term from the original test and
holds for every pair of relations in the representation algebra.
% END Q70C2

% BEGIN Q70C3
\paragraph{Q70C3. A relation in one corner separates the two cyclotomic weights.}
Suppose both blocks are nonempty, and choose $g_0\in G$ with
$\chi(g_0)\ne1$, as required by Choice. In either off-diagonal
corner fix $i,j$. If the $ij$ entry of
$D=\sum_d(\alpha_d K_d+\beta_d L_d)$ vanishes on $I$, the
entire chosen corner vanishes there. Indeed
$D(puq)=\pi(p)D(u)\pi(q)$ for $u\in I$. Left and right lifts
of matrix units move any entry in that corner to $ij$.
The corner of $D$ therefore descends to a derivation
$B_E\to V_{12}$ or $B_E\to V_{21}$. Such a derivation is
inner. For the central idempotent $P$ acting as identity on the
left and zero on the right of that corner, differentiate
$P b=b P$. This gives $D(b)=bD(P)-D(P)b$.
After multiplying on the right by $r(h)^{-1}$, the restricted
corner of $\sum\alpha_d\kappa_d+\sum\beta_d\lambda_d$ is a
coboundary on $H$ for the untwisted action.

In the space of $H$-cocycles with values in this corner modulo
coboundaries let $T_g$ send a representative $u(h)$ to
$r(g)u(g^{-1}hg)r(g)^{-1}$. It preserves coboundaries. If $u$
comes from a global cocycle for $\theta\operatorname{Ad}r$,
where $\theta|_H=1$, expansion of $hg=g(g^{-1}hg)$ shows
\[
 \theta(g)(T_g u)(h)=u(h)+r(h)u(g)r(h)^{-1}-u(g).
\]
The two summands in our relation therefore have classes in the
$1$ and $\chi(g_0)^{-1}$ eigenspaces of $T_{g_0}$.
Their sum vanishes, so both vanish. Each corner is a global
coboundary in its respective twist. Subtract its $H$-primitive
as a global coboundary, and the remaining cocycle vanishes on
$H$. Comparing $hg$ with $g(g^{-1}hg)$ shows each of its values
is $H$-invariant. Such a corner commutes with $B_E$ by the span
condition. Commuting with $P$ forces it to vanish.

Write $d_\theta Y(g)=\theta(g)r(g)Yr(g)^{-1}-Y$. The form gives
$(d_\theta Y(\sigma g))^T=d_\theta(Y^T)(g)$ for either twist. A primitive in
one corner extends to both as $Y+\delta Y^T$ for the prescribed
transpose sign $\delta$ of that cocycle. Thus $\sum\alpha_d\kappa_d$
has a skew odd primitive and $\sum\beta_d\lambda_d$ has an odd
primitive with sign $s$. Independence in Q70C1 forces
$\alpha_d=0$ for every $d$.
% END Q70C3

% BEGIN Q70C4
\paragraph{Q70C4. Detection for the original twisted cocycles.}
For fixed $v\in I$ put $\alpha_d=L_d(v)_{ji}$ and
$\beta_d=K_d(v)_{ji}$ in Q70C3. The identity in Q70C2 applies
to every $u$, hence $L_d(v)_{ji}=0$ for all $v,d$. Now apply
Q70C3 in the reverse corner with only the coefficient of $L_d$
nonzero and equal to one. It follows that every $\lambda_d$ is
$d_\chi Y_d$ with $Y_d$ odd and $sY_d^T=Y_d$.
When $a=n$ every odd component is already zero and $Y_d=0$ works.

Put $Y'=\sum f_dY_d$ and subtract the coboundary
$\chi(g)R(g)Y'R(g)^{-1}-Y'$ from $Z$. The remaining cocycle is
$\eta z(g)$ with $z(g)\in B_E$. It has transpose sign $s$ by
the form for $R$. Its trace test on $H$ still vanishes, since
the removed term multiplied on the right by $R(h)$ is a
commutator. Thus $z$ is a $\chi\operatorname{Ad}r$-cocycle with
$\operatorname{tr}(z(h)r(h))=0$ for all $h\in H$.

On $E[H]$ the linear map $D(h)=z(h)r(h)$ satisfies the derivation
rule of Q70C2. For $u\in I$ and $v\in E[H]$,
$0=\operatorname{tr}D(uv)=\operatorname{tr}(D(u)\pi(v))$.
The trace pairing on $B_E$ is nondegenerate, by testing against
matrix units within each block. Hence $D$ descends to $B_E$.
A derivation of this product annihilates its central block idempotents since
$D(P)=D(P^2)=2P D(P)$, which has zero components on both
blocks. On each matrix factor it is inner by the matrix-unit
calculation in CCH4. Subtracting the corresponding even coboundary
from $z$ leaves a cocycle vanishing on $H$. As in Q70C3 its values
centralize $B_E$. They are scalar on each block, so conjugation
by $r(G)$ acts trivially on them. This residual cocycle factors
through the abelian group $G/H$. Comparing the two orders of
multiplication with $g_0$ shows it is the coboundary of its
value at $g_0$ divided by $\chi(g_0)-1$. Consequently
$z=d_\chi Y_0$ for an even primitive, and replacing $Y_0$ by
$(Y_0+sY_0^T)/2$ gives the required symmetry.

The full primitive $Y'+\eta Y_0$ lies in the original dual
matrix space and has transpose sign $s$. For an integral cocycle
one power of $t$ clears all its finitely many coefficient
denominators, giving $t^b Z\in\mathrm{TwB}(k_x,A)$.
For any $U\subset H$ dense in $H$, vanishing of the detector on
$U$ implies vanishing on $H$ because its trace coefficient is continuous
modulo every $t^j$, and $A$ is separated. This proves the native
$\mathrm{Detection}$, with no finiteness or separability input.
% END Q70C4

% BEGIN Q70C5
\paragraph{Q70C5. A common finite detector in the simultaneous separable locus.}
Now take a family $x_i$ indexed by $0\le i\le h$, with all the
hypotheses of Choice. Let $U$ consist of the $g\in H$ for which
all the characteristic polynomials of $r_i(g)$ are separable over
their respective $E_i$. Each separate condition is open in $H$ since it is the nonvanishing of the discriminant, a polynomial in
the entries, hence it is witnessed modulo some $t_i^N$.
Each condition is dense by $x_i.N3$. Their finite intersection
is dense. Starting with a nonempty relatively open set, intersect
successively with these open dense sets. Choose $k_i$ by Q70C1.
Then $U$ detects every $k_i$ by Q70C4.

At each point the algebra with basis $1,e_d,f_d,\eta$ is finite
free over $A_i$. Apply the finite-generation argument of Q70C1 to
its full matrix space with action $\chi\operatorname{Ad}R_i$.
This action is continuous, since $r_i$ and all $\kappa_{i,d}$
are continuous and $R_i(g)^{-1}=R_i(g^{-1})$.
Take a finite list of module generators for the resulting
cocycles, without imposing the polarization condition. For
$u\in U$, take the row of values of
$Z\mapsto[\eta]\operatorname{tr}(Z(u)R_i(u))$ on that list.
Choose a basis of the span of the rows over $E_i$ from among
the rows themselves. Vanishing on this finite selection implies
vanishing on all of $U$, by linearity, also for every actual
twisted cocycle. Take the union of the selections over all
points, enumerate it as $\gamma_1,\ldots,\gamma_q$, and keep
this list fixed. Its elements belong to $H$ and every point's
selected tests occur, so it detects all $k_i$. The construction
includes the case $q=0$.
% END Q70C5

% BEGIN Q70C6
\paragraph{Q70C6. Ramified coefficient extension of the odd generators and detector.}
Fix one point, and a finite injective local map $j:A\to A'$
between equal-characteristic power-series DVRs with finite
residue fields of characteristic $\ell$. Apply $j$ to all entries
of $r,w_v,w_v^{-1},\psi_v$ and keep the block size and column
types. The resulting data give $x':\mathrm{NP}(F,n)$.
In fact $j(t)=u(t')^e$ for $e\ge1$ and a unit $u$. Continuity
passes to the new data by this identity. Applying the old span
test to the block matrix units and transporting the expressions
gives the span test over $A'$ after multiplying by a common
power of $t'$. All form, inertia, character, column support
and triangularity identities transport directly. The cyclotomic
scalars and $s$ agree under $j$ by the common prime-field
factorization, and injectivity preserves both character exclusions.
Hence this gives the original $\mathrm{Ext}$ fields.

Take the generator cocycles and their local witnesses to be
the entrywise images of $\kappa_d,Y_{d,v}$. To verify the other
fields, use an $A$-basis of $A'$, which exists because $A'$ is a
finite torsion-free module over a DVR. Each coordinate projection
is $A$-linear and continuous since $(t')^{eN}A'=t^NA'$.
Acting matrices, flag matrices, transpose and scalars are
defined over $A$, so every defining linear equation can be
projected coefficient by coefficient. In particular project any
local witness together with its cocycle. An integral relation
among the new generators modulo a skew odd coboundary yields
a relation downstairs for each coordinate, making all the
coefficients zero. Conversely, project each new Selmer cocycle
to the old Selmer module, apply the old generation field to
every coordinate, use the largest exponent, and recombine the
coefficients and skew primitives. This gives generation after
multiplication by $j(t)^N$. Dividing by the unit $u^N$ gives
the required power $(t')^{eN}$. Thus the transported data form
a genuine $k':\mathrm{Kap}(x')$.

For the same reason, applying the coordinate projections to every
basis coefficient $1,e_d,f_d,\eta$ sends an integral twisted
cocycle upstairs to one downstairs. Coefficients of $R$, its
inverse and all the multiplications in its action descend, and the
filtration comparison gives the prescribed adic continuity for
each projection. Projection also commutes with the trace functional
since the matrix multiplied on its right descends. If the tests
upstairs vanish on the list, the old detector supplies a primitive
after a power of $t$ for every projected cocycle. Multiply up to
the largest exponent and recombine the primitives. Their signs and
block supports are preserved, and the same unit adjustment gives
the original $\mathrm{Detects}(k',\{\gamma_j\})$ over $A'$.
% END Q70C6

% BEGIN Q70C7
\paragraph{Q70C7. Common residual cardinality and the original Choice output.}
For each $i$ take the compositum of the splitting fields of the
characteristic polynomials at the fixed $q$ elements. These are
finite separable extensions of $E_i$, by Q70C5. The integral
closure of $A_i$ in this field is a finite complete DVR by CCH9.
Each root belongs to it by monicity, and the roots of each
polynomial are pairwise distinct. By CCH10 this DVR is a power
series ring over its residue field. Lift the finite field of order
$\ell^{d_i}$ along the simple roots of $X^{\ell^{d_i}}-X$ and
expand in a uniformizer. In particular $d_i>0$, also for the
empty splitting list. Take a common positive multiple $D$ of all
the $d_i$. Extend the constant field at point $i$ to
$\mathbb F_{\ell^D}$, keeping its new uniformizer. Composing the
two maps gives the finite injective local map used in Q70C6.

Label the roots of each polynomial in any order by
$\mathrm{Fin}(n)$. Under this injection they remain distinct,
and the factorization into the listed linear factors remains the
characteristic polynomial of the transported matrix by functoriality.
Every point now has residue cardinality $\ell^D$. Together with
the one list in $H$, the transported packages and their detectors
from Q70C5--Q70C6 give every output of $\mathrm{Pa.Choice}$,
for every original $h,n,F$ and all permissible block sizes.
N69E4 can use this result for its whole chain. The two arithmetic
bounds and the support theorem $\mathrm{Raw}$ in that implication
are separate inputs.
% END Q70C7
